\newif\ifcircuitsflip

\makeatletter

% drawing styles
\tikzset{
    circuits/general/.style = {thick},
    wire/.style={circuits/general},
    cable/.style={circuits/general, cap = round},
    %rfcable/.style={circuits/general, double, double distance = 0.6pt, cap = round},
    rfcable/.style={circuits/general, double, cap = round},
    arrowwire/.style = {
        circuits/general,
        ->,
        >=stealth,
    },
    lwire/.style={circuits/general, line cap = rect},
    scriptsize/.style={font=\scriptsize},
    ---/.style args={#1 and #2}{to path={-- ++(#1, 0) |- ($(\tikztotarget)+(#2, 0)$)}},
    ---/.default = 0cm and 0cm,
    -|-/.style args={#1 and #2}{to path={-| ([xshift=#1, yshift=#2]$(\tikztostart)!0.5!(\tikztotarget)$) |- (\tikztotarget) \tikztonodes}},
    -|-/.default = 0cm and 0cm,
    |-|/.style args={#1 and #2}{to path={|- ([xshift=#1, yshift=#2]$(\tikztostart)!0.5!(\tikztotarget)$) -| (\tikztotarget) \tikztonodes}},
    |-|/.default = 0cm and 0cm,
    angleconnect/.style args={#1 and #2}{to path={-- ++(#1, 0) -- ($(\tikztotarget)-(#2, 0)$) -- (\tikztotarget)}},
    angleconnect/.value required = true,
    angleconnectx/.style args={#1 and #2}{to path={-- ++(#1, 0) -- ($(\tikztotarget)-(#2, 0)$) -- (\tikztotarget)}},
    angleconnectx/.value required = true,
    angleconnecty/.style args={#1 and #2}{to path={-- ++(0, #1) -- ($(\tikztotarget)-(0, #2)$) -- (\tikztotarget)}},
    angleconnecty/.value required = true,
    tmodule/.style args ={#1 and #2}{draw, rectangle, wire, minimum width = #1, minimum height = #2},
    tmodule/.default = {0cm and 0cm},
    module width/.store in=\circuits@module@width,
    module height/.store in=\circuits@module@height,
    module width = 1cm,
    module height = 1cm,
    module/.style = {
        wire,
        draw, rectangle,
        align = center,
        minimum width = \circuits@module@width,
        minimum height = \circuits@module@height
    },
    current arrow/.style = {
        postaction = decorate,
        decoration = {
            markings,
            mark = at position #1 with {\arrow{stealth}}
        }
    },
    current arrow/.default = 0.5,
    current arrow reversed/.style = {
        postaction = decorate,
        decoration = {
            markings,
            mark = at position #1 with {\arrowreversed{stealth}}
        }
    },
    current arrow reversed/.default = 0.5,
    bit line/.style = {
        wire,
        postaction = decorate,
        decoration = {
            markings,
            mark = at position 0.5 with { \draw (-#1, -#1) -- (#1, #1); }
        }
    },
    bit line/.default = 2pt,
    % flip components (effect depends on shape)
    flip/.is if = circuitsflip,
    xmirror/.style = {xscale = -1},
    ymirror/.style = {yscale = -1}
}

\pgfkeys{
    % general circuits options
    /tikz/circuits/line width/.initial = 0.8pt,
    /tikz/circuits/dot/radius/.initial = 0.05cm,
    /tikz/dot/.style = {dotshape},
    /tikz/circuits/port/radius/.initial = 0.05cm,
    %/tikz/port/.style = {portshape}
    /tikz/port/.style = {wire, draw=black, fill=white, circle, inner sep = 0cm, minimum size = 0.1cm}
}

\pgfdeclareshape{dotshape}
{
    \saveddimen{\radius}{\pgf@x=\pgfkeysvalueof{/tikz/circuits/dot/radius}}
    \anchor{center}{\pgfpointorigin}
    \anchor{north}{\pgfpointorigin}
    \anchor{south}{\pgfpointorigin}
    \anchor{west}{\pgfpointorigin}
    \anchor{east}{\pgfpointorigin}
    \anchor{north east}{\pgfpointorigin}
    \anchor{north west}{\pgfpointorigin}
    \anchor{south east}{\pgfpointorigin}
    \anchor{south west}{\pgfpointorigin}
    \backgroundpath{
        \pgfpathcircle{\pgfpointorigin}{\radius}
        \pgfusepath{fill}
    }
}
\pgfdeclareshape{portshape}
{
    \saveddimen{\radius}{\pgf@x=\pgfkeysvalueof{/tikz/circuits/port/radius}}
    \anchor{center}{\pgfpointorigin}
    \backgroundpath{
        \color{white}
        \pgfsetstrokecolor{black}
        \pgfsetlinewidth{\pgfkeysvalueof{/tikz/circuits/line width}}
        \pgfpathcircle{\pgfpointorigin}{\radius}
        \pgfusepath{fill, stroke}
    }
}

\def\circuitskeysvalueof#1{\pgfkeysvalueof{/tikz/circuits/.cd, #1}}

\pgfkeys{
    /tikz/circuits/scale/.initial = 1
}

\input{pkg/symbols/ota}
\input{pkg/symbols/opamp}
\input{pkg/symbols/comparator}
\input{pkg/symbols/capacitor}
\input{pkg/symbols/resistor}
\input{pkg/symbols/inductor}
\input{pkg/symbols/mosfet}
\input{pkg/symbols/current_source}
\input{pkg/symbols/ground}
\input{pkg/symbols/vdd}
\input{pkg/symbols/impedance}
\input{pkg/symbols/amplifier}
\input{pkg/symbols/tgate}
\input{pkg/symbols/mixer}
\input{pkg/symbols/oscillator}
\input{pkg/symbols/antenna}
\input{pkg/symbols/diode}
\input{pkg/symbols/gyrator}
\input{pkg/symbols/inverter}
\input{pkg/symbols/switch}
\input{pkg/symbols/magglass}
\input{pkg/symbols/balun}
\input{pkg/symbols/chip}
\input{pkg/symbols/probe}
\input{pkg/symbols/instrument}
\input{pkg/symbols/supply}
\input{pkg/symbols/computer}
\input{pkg/symbols/pcb}
\input{pkg/symbols/adc}
\input{pkg/symbols/dac}
\input{pkg/symbols/sumpoint}
\input{pkg/symbols/logic_gates}
\newif\if@circuits@sequentiallogic@drawreset
\tikzset{
    circuits/sequential logic/draw reset/.is if=@circuits@sequentiallogic@drawreset,
}
\input{pkg/symbols/dflipflop}
\input{pkg/symbols/counter}
\input{pkg/symbols/register}
\input{pkg/symbols/statemachine}

% vim: ft=plaintex nowrap

\input{pkg/symbols/mux}
\makeatother

% vim: ft=plaintex nowrap
